%
% teil3.tex -- Eigenschaften der Bessel-Funktionen
%
% (c) 2026 Gian Cavegn, OST Ostschweizer Fachhochschule
%
% !TEX root = ../../buch.tex
% !TEX encoding = UTF-8
% !TeX spellcheck = de_CH
%
\section{Eigenschaften der Bessel-Funktionen
\label{bessel:section:eigenschaften}}
\kopfrechts{Eigenschaften}
%
% Dieser Abschnitt ist der Werkzeugkasten. Alles, was hier steht, wird
% später gebraucht.
%   - Rekursionsformeln       -> Rechnen allgemein
%   - Integraldarstellung     -> teil6 (Kepler), das ist die Brücke!
%   - Nullstellen             -> Membran, teil5 (Airy)
%   - Orthogonalität          -> Fourier-Bessel-Reihe, teil6
%   - Asymptotik              -> teil5 und Vergleich mit cos
%
Die Reihe~\eqref{bessel:loesung:eqn:jnu} definiert die Bessel-Funktionen, ist zum Rechnen aber unhandlich.
Dieser Abschnitt stellt die Eigenschaften zusammen, mit denen man tatsächlich arbeitet.
Es lohnt sich, dabei immer den Vergleich mit den trigonometrischen Funktionen im Auge zu behalten.
Fast jede Eigenschaft von $J_\nu$ hat ein Gegenstück bei $\cos$ und $\sin$.
%
% Rekursionsformeln
%
\subsection{Rekursionsformeln
\label{bessel:eigenschaften:subsection:rekursion}}
Aus der Reihendarstellung lassen sich zwei Ableitungsformeln direkt ablesen.

\begin{satz}
\label{bessel:eigenschaften:satz:ableitung}
Für alle $\nu$ gilt
\begin{align}
\frac{d}{dx}
\bigl( x^{\nu} J_\nu(x) \bigr)
&=
x^{\nu} J_{\nu-1}(x),
\label{bessel:eigenschaften:eqn:ableitung1}
\\
\frac{d}{dx}
\bigl( x^{-\nu} J_\nu(x) \bigr)
&=
-x^{-\nu} J_{\nu+1}(x).
\label{bessel:eigenschaften:eqn:ableitung2}
\end{align}
\end{satz}
 
\begin{proof}
Aus~\eqref{bessel:loesung:eqn:jnu} folgt
\[
x^\nu J_\nu(x)
=
\sum_{k=0}^\infty
\frac{(-1)^k}{k!\,\Gamma(\nu+k+1)}
\frac{x^{2k+2\nu}}{2^{2k+\nu}}.
\]
Gliedweises Differenzieren ist innerhalb des Konvergenzbereiches erlaubt und ergibt
\[
\frac{d}{dx}
\bigl( x^\nu J_\nu(x) \bigr)
=
\sum_{k=0}^\infty
\frac{(-1)^k (2k+2\nu)}{k!\,\Gamma(\nu+k+1)}
\frac{x^{2k+2\nu-1}}{2^{2k+\nu}}.
\]
Mit $\Gamma(\nu+k+1)=(\nu+k)\Gamma(\nu+k)$ kürzt sich der Faktor $2(k+\nu)$ weg, und es bleibt genau die Reihe von $x^\nu J_{\nu-1}(x)$.
Die zweite Formel folgt analog.
\end{proof}
 
Führt man die Produktregel aus und eliminiert einmal die Ableitung und einmal $J_\nu$, erhält man die beiden für die Praxis wichtigsten Beziehungen.
 
\begin{korollar}
\label{bessel:eigenschaften:korollar:rekursion}
Es gilt
\begin{align}
J_{\nu-1}(x) + J_{\nu+1}(x)
&=
\frac{2\nu}{x} J_\nu(x),
\label{bessel:eigenschaften:eqn:dreiterm}
\\
J_{\nu-1}(x) - J_{\nu+1}(x)
&=
2 J_\nu'(x).
\label{bessel:eigenschaften:eqn:ableitungrek}
\end{align}
\end{korollar}
 
Die erste Formel ist eine Dreitermrekursion.
Kennt man $J_0$ und $J_1$, so lassen sich alle weiteren Ordnungen daraus berechnen.
Aus~\eqref{bessel:eigenschaften:eqn:ableitungrek} folgt für $\nu=0$ wegen $J_{-1}=-J_1$ der oft gebrauchte Spezialfall
\begin{equation}
J_0'(x) = -J_1(x),
\label{bessel:eigenschaften:eqn:j0strich}
\end{equation}
das genaue Gegenstück zu $\cos' = -\sin$.
 
%
% Erzeugende Funktion und Integraldarstellung
%
\subsection{Erzeugende Funktion und Integraldarstellung
\label{bessel:eigenschaften:subsection:erzeugend}}
Für ganzzahlige Ordnungen gibt es eine überraschend kompakte Beschreibung aller $J_n$ auf einmal.
 
\begin{satz}
\label{bessel:eigenschaften:satz:erzeugend}
Für $t\ne 0$ gilt
\begin{equation}
e^{\frac{x}{2}\left(t-\frac{1}{t}\right)}
=
\sum_{n=-\infty}^{\infty}
J_n(x)\, t^n.
\label{bessel:eigenschaften:eqn:erzeugend}
\end{equation}
\end{satz}
 
\begin{proof}
Man entwickelt die beiden Faktoren $e^{xt/2}$ und $e^{-x/(2t)}$ je in ihre Exponentialreihe, multipliziert die Reihen aus und sammelt die Terme mit gleicher Potenz von $t$.
Der Koeffizient von $t^n$ ist dann genau die Reihe
\eqref{bessel:loesung:eqn:jnu}.
\end{proof}
% [TODO: das Cauchy-Produkt explizit ausschreiben, mindestens für n>=0]
 
Setzt man $t=e^{i\tau}$, so wird $t-1/t = 2i\sin\tau$, und \eqref{bessel:eigenschaften:eqn:erzeugend} geht über in
\begin{equation}
e^{ix\sin\tau}
=
\sum_{n=-\infty}^{\infty}
J_n(x)\, e^{in\tau}.
\label{bessel:eigenschaften:eqn:jacobianger}
\end{equation}
Das ist die \emph{Jacobi-Anger-Entwicklung}.
\index{Jacobi-Anger-Entwicklung}%
Die rechte Seite ist eine Fourier-Reihe in $\tau$, und die $J_n(x)$ sind ihre Fourier-Koeffizienten.
Die übliche Formel für Fourier-Koeffizienten liefert damit sofort
 
\begin{korollar}
\label{bessel:eigenschaften:korollar:integral}
Für $n\in\mathbb{Z}$ gilt
\begin{equation}
J_n(x)
=
\frac{1}{2\pi}
\int_{-\pi}^{\pi}
e^{i(x\sin\tau - n\tau)}\,d\tau
=
\frac{1}{\pi}
\int_0^{\pi}
\cos(n\tau - x\sin\tau)\,d\tau.
\label{bessel:eigenschaften:eqn:integral}
\end{equation}
\end{korollar}
 
Diese Integraldarstellung ist historisch die ursprüngliche.
Bessel hat die Funktionen in dieser Form eingeführt, und sie ist die Brücke zur Kepler-Gleichung in Abschnitt~\ref{bessel:section:kepler}.
Der Ausdruck $n\tau - x\sin\tau$ im Integranden ist bereits die linke Seite der Kepler-Gleichung.
 
%
% Nullstellen
%
\subsection{Nullstellen
\label{bessel:eigenschaften:subsection:nullstellen}}
Die Bessel-Funktionen oszillieren, und ihre Nullstellen bestimmen in jeder Anwendung die möglichen Eigenfrequenzen oder Beugungsminima.
 
\begin{satz}
\label{bessel:eigenschaften:satz:nullstellen}
Für $\nu\ge 0$ hat $J_\nu$ unendlich viele positive Nullstellen
\[
0 < j_{\nu,1} < j_{\nu,2} < j_{\nu,3} < \dots,
\]
und es gilt $j_{\nu,m}\to\infty$.
Zwischen je zwei aufeinanderfolgenden Nullstellen von $J_\nu$ liegt genau eine Nullstelle von $J_{\nu+1}$ und umgekehrt.
\end{satz}
% [TODO: Beweis über die Sturmsche Vergleichstheorie oder über die
%        Substitution u = sqrt(x) y, die auf u'' + (1 + (1/4-nu^2)/x^2)u = 0
%        führt. Diese Substitution ist auch für die Asymptotik nützlich.]
 
Die Nullstellen sind nicht in geschlossener Form angebbar, sie müssen numerisch bestimmt werden.
Die ersten Werte sind
\begin{center}
\begin{tabular}{c|ccc}
$\nu$ & $j_{\nu,1}$ & $j_{\nu,2}$ & $j_{\nu,3}$ \\
\hline
$0$ & $2.4048$ & $5.5201$ & $8.6537$ \\
$1$ & $3.8317$ & $7.0156$ & $10.1735$ \\
$2$ & $5.1356$ & $8.4172$ & $11.6198$
\end{tabular}
\end{center}
% [TODO: Tabelle als richtige Table-Umgebung mit Caption und Label,
%        Quelle für die Werte angeben (Abramowitz-Stegun)]
 
Anders als bei $\sin$ sind die Nullstellen nicht äquidistant, ihre Abstände nähern sich aber für grosse $x$ dem Wert $\pi$ an.
Das ist die Erklärung für den unharmonischen Klang einer Trommel.
Die Eigenfrequenzen der kreisförmigen Membran vom Radius $a$ ergeben sich aus der Randbedingung $R(a)=0$, also aus $J_n(ka)=0$, zu
\begin{equation}
\omega_{nm}
=
\frac{c}{a}\, j_{n,m}.
\label{bessel:eigenschaften:eqn:eigenfrequenzen}
\end{equation}
Das Verhältnis der beiden tiefsten rotationssymmetrischen Moden ist $j_{0,2}/j_{0,1}\approx 2.295$ und damit kein rationales Vielfaches des Grundtons.
Bei der Saite dagegen liegen die Nullstellen des Sinus äquidistant, und die Obertöne sind ganzzahlige Vielfache der Grundfrequenz.
% [TODO: Abbildung der ersten Membranmoden (0,1), (1,1), (2,1), (0,2)
%        mit Knotenlinien]
 
%
% Orthogonalität
%
\subsection{Orthogonalität
\label{bessel:eigenschaften:subsection:orthogonalitaet}}
Wie die trigonometrischen Funktionen bilden auch die Bessel-Funktionen ein Orthogonalsystem, allerdings mit einer Gewichtsfunktion.
 
\begin{satz}
\label{bessel:eigenschaften:satz:orthogonalitaet}
Sind $j_{\nu,m}$ und $j_{\nu,l}$ Nullstellen von $J_\nu$, so gilt
\begin{equation}
\int_0^1
J_\nu(j_{\nu,m}s)\,
J_\nu(j_{\nu,l}s)\,
s\,ds
=
\begin{cases}
0 & \text{für } m\ne l,\\[3pt]
\dfrac{1}{2}\,J_{\nu+1}(j_{\nu,m})^2 & \text{für } m=l.
\end{cases}
\label{bessel:eigenschaften:eqn:orthogonalitaet}
\end{equation}
\end{satz}
% [TODO: Beweis über die Lagrange-Identität / partielle Integration
%        der Bessel-Gleichung in selbstadjungierter Form
%        (x R')' + (k^2 x - nu^2/x) R = 0]
 
Das Gewicht $s$ im Integral ist kein Schönheitsfehler, sondern das Flächenelement der Polarkoordinaten, $dA = r\,dr\,d\varphi$.
Die Orthogonalität ist also die natürliche Orthogonalität auf der Kreisscheibe.
 
Daraus ergibt sich die \emph{Fourier-Bessel-Reihe}.
\index{Fourier-Bessel-Reihe}%
Eine auf $[0,1]$ hinreichend gutartige Funktion $f$ lässt sich als
\begin{equation}
f(s)
=
\sum_{m=1}^{\infty}
c_m J_\nu(j_{\nu,m}s),
\qquad
c_m
=
\frac{2}{J_{\nu+1}(j_{\nu,m})^2}
\int_0^1 f(s)\,J_\nu(j_{\nu,m}s)\,s\,ds,
\label{bessel:eigenschaften:eqn:fourierbessel}
\end{equation}
entwickeln.
Damit lassen sich Anfangsbedingungen einer kreisförmigen Membran nach Eigenmoden zerlegen, genau wie eine Fourier-Reihe das für die Saite leistet.
 
%
% Asymptotisches Verhalten
%
\subsection{Asymptotisches Verhalten
\label{bessel:eigenschaften:subsection:asymptotik}}
Für grosse Argumente werden die Bessel-Funktionen fast zu trigonometrischen Funktionen.
 
\begin{satz}
\label{bessel:eigenschaften:satz:asymptotik}
Für $x\to\infty$ gilt
\begin{align}
J_\nu(x)
&=
\sqrt{\frac{2}{\pi x}}
\cos\Bigl(x-\frac{\nu\pi}{2}-\frac{\pi}{4}\Bigr)
+
O(x^{-3/2}),
\label{bessel:eigenschaften:eqn:asymptotikj}
\\
Y_\nu(x)
&=
\sqrt{\frac{2}{\pi x}}
\sin\Bigl(x-\frac{\nu\pi}{2}-\frac{\pi}{4}\Bigr)
+
O(x^{-3/2}).
\label{bessel:eigenschaften:eqn:asymptotiky}
\end{align}
\end{satz}
% [TODO: Herleitung mindestens plausibel machen über die Substitution
%        u(x) = sqrt(x) y(x), die auf u'' + (1 + (1/4-nu^2)/x^2) u = 0
%        führt; für grosse x ist das die Schwingungsgleichung u'' + u = 0.]
 
Die Substitution $u(x)=\sqrt{x}\,y(x)$ macht diese Aussage verständlich.
Sie überführt die Bessel-Gleichung in
\[
u''
+
\Bigl(
1 + \frac{\tfrac14-\nu^2}{x^2}
\Bigr) u
=
0,
\]
und für grosse $x$ ist das bis auf einen verschwindenden Term die gewöhnliche Schwingungsgleichung $u''+u=0$.
Die Lösung ist also asymptotisch eine Schwingung, und der Faktor $1/\sqrt{x}$ kommt von der Rücksubstitution.
 
Die Abklingrate $1/\sqrt{x}$ ist physikalisch einleuchtend.
Eine von einem Punkt ausgehende Kreiswelle verteilt ihre Energie auf einen Kreis mit Umfang $2\pi r$, die Energiedichte fällt also wie $1/r$ und die Amplitude wie $1/\sqrt{r}$.
Wirft man einen Stein ins Wasser, sieht man genau das.
 
Zusammen mit~\eqref{bessel:loesung:eqn:kleinx} ist das Verhalten von $J_\nu$ damit in beiden Grenzfällen bekannt.
Nahe null verhält es sich wie eine Potenz $x^\nu$, weit draussen wie eine gedämpfte Kosinusschwingung.
Der Bereich dazwischen ist der interessante, und dort hilft nur die Reihe oder die Numerik.